\subsection{Orthogonal projectors and involutions.}

Throughout this no., one assumes that the scalar 2 is \(invertible\) in A (for example that A is a field of characteristic \(\neq 2\) ), and that \(\Phi\) is a Hermitian form \(nondegenerate\) on E. One denotes by \(\frac{1}{2}\) the inverse of 2.

Lemma 2. --- For an endomorphism u of E to be such that \(u^{2}=1\) , it is necessary and sufficient that \(\frac{1}{2}(1-u)\) be a projector in E; then u is the difference of the two projectors \(\frac{1}{2}(1+u)\) and \(\frac{1}{2}(1-u)\) .

Indeed, in the ring \(\mathfrak{L}(\mathrm{E})\) , the relation \(\left(\frac{1}{2}(1 - u)\right)^2 = \frac{1}{2}(1 - u)\) is equivalent to \(u^2 = 1\) . The rest is trivial.

An endomorphism \(u\) of E such that \(u^2 = 1\) (which is then necessarily an automorphism of E equal to its inverse) is called an involution. Set \(v = \frac{1}{2}(1 - u)\) , \(U^- = v(E)\) , \(U^+ = v^{-1}(0)\) ( \(= w(E)\) by setting \(w = \frac{1}{2}(1 + u)\) ); one knows that E is the direct sum of \(U^+\) and of \(U^-\) (chap. VIII, § 1, no. 1), and one has \(u(x) = x\) in \(U^+\) , \(u(x) = -x\) in \(U^-\) . When A is a field and E is of finite dimension, it follows, since A is of characteristic \(\neq 2\) , that the only eigenvectors \(\neq 0\) of \(u\) are the elements ≠ 0 in U⁺ or in U⁻; they correspond respectively to the eigenvalues +1 and −1.

PROPOSITION 4. --- Let \(u \in \mathbf{GL}(\mathrm{E})\) be an involution. The following properties are equivalent:

\begin{enumerate}
\def\labelenumi{\alph{enumi})}
\item
  u belongs to the unitary group associated with \(\Phi\) ;
\item
  the submodules \(U^{+} = \frac{1}{2}(1 + u)(E)\) and \(U^{-} = \frac{1}{2}(1 - u)(E)\) are orthogonal (and consequently non-isotropic).
\end{enumerate}

Moreover, if A is a field and E is finite-dimensional, properties a) and b) are equivalent to:

\begin{enumerate}
\def\labelenumi{\alph{enumi})}
\setcounter{enumi}{2}
\tightlist
\item
  \(u = u^{*}\) .
\end{enumerate}

Indeed, for \(x \in \mathrm{U}^{+}\) and \(y \in \mathrm{U}^{-}\) , the relation \(\Phi(u(x), u(y)) = \Phi(x, y)\) gives \(2\Phi(x, y) = 0\) , hence \(a)\) implies \(b)\) . Conversely, we obviously have \(\Phi(u(x), u(y)) = \Phi(x, y)\) when \(x\) and \(y\) are both in \(\mathrm{U}^{+}\) or both in \(\mathrm{U}^{-}\) , and, given \(b)\) , this relation is still true when one of them is in \(\mathrm{U}^{+}\) and the other in \(\mathrm{U}^{-}\) ; since E is the direct sum of \(\mathrm{U}^{+}\) and \(\mathrm{U}^{-}\) , we see that \(b)\) implies \(a)\) . Finally, when E is a finite-dimensional vector space, the adjoint \(u^{*}\) is defined since \(\Phi\) is nondegenerate; the relation \(a)\) is equivalent to \(uu^{*} = 1\) ( \(§ 1, n^{\circ} 8\) , cor. of prop. 8); since \(u^{2} = 1\) by hypothesis, \(a)\) and \(c)\) are equivalent.

COROLLARY 1. --- Suppose that A is a field and that E is finite-dimensional. The map \(u \to \frac{1}{2}(1 + u)(E)\) is a bijection from the set of involutions u belonging to the unitary group associated with \(\Phi\) onto the set of non-isotropic subspaces of E; the subspace U \(^+\) corresponding to u is the set of elements of E invariant under u.

By prop. 4, it suffices to show that every non-isotropic subspace M of E is the set of vectors invariant under an involution \(u \in \mathbf{U}(\Phi)\) , and that this involution is unique. Now E is the direct sum of M and of \(\mathbf{M}^0\) ( \(\S 4\) , no. 1, cor. of prop. 1), and one necessarily has \(u(x) = x\) for \(x \in \mathbf{M}\) and \(u(x) = -x\) for \(x \in \mathbf{M}^0\) by virtue of prop. 4; these relations determine \(u\) uniquely, and the endomorphism \(u\) thus determined evidently answers the question (prop. 4).

One says that the involution \(u\) thus determined is the symmetry with respect to the non-isotropic subspace M.

COROLLARY 2. --- Suppose that A is a field and that E is of finite dimension. For a projector v in E to be such that v(E) and \(v^{-1}(0)\) are orthogonal (and consequently non-isotropic), it is necessary and sufficient that v = v*.

It suffices to apply prop. 4 to the involution \(u = 1 - 2v\) .

A projector satisfying the condition of corollary 2 is called an orthogonal projector for \(\Phi\) .

